% Generated draft: requires mathematical and source review.
\documentclass[11pt]{article}
\usepackage[margin=25mm]{geometry}
\usepackage{fontspec}
\setmainfont{Latin Modern Roman}
\usepackage{amsmath,amssymb,mathtools,unicode-math}
\setmathfont{Latin Modern Math}
\usepackage{enumitem,needspace,xurl,hyperref}
\setlength{\parindent}{0pt}
\setlength{\parskip}{5pt}
\setlength{\emergencystretch}{3em}
\newcommand{\N}{\mathbb{N}}
\newcommand{\Z}{\mathbb{Z}}
\newcommand{\Q}{\mathbb{Q}}
\newcommand{\R}{\mathbb{R}}
\newcommand{\C}{\mathbb{C}}
\newcommand{\F}{\mathbb{F}}
\newcommand{\A}{\mathbb{A}}
\newcommand{\PP}{\mathbb{P}}
\newcommand{\E}{\mathbb{E}}
\newcommand{\eps}{\varepsilon}
\newcommand{\dd}{\,\mathrm{d}}
\begin{document}
\section*{UnsolvedMath ID: 30004836}
\textbf{Source problem code:} catalogue entry
\subsection{Definitions and mathematical statement}
The supplied record concerns Nichols algebras and a sheaf-theoretic correspondence, but does not specify the underlying field, braided category, geometric spaces, or correspondence. The symbols \(F\), \(K_n\), \(V_n\), and \(T_{!*}\) are retained as source notation: their definitions are absent from the supplied record. In particular, no interpretation of \(T_{!*}\) as a particular extension functor is assumed.

For terminology only, if \(V\) is a braided vector space over a field \(k\) in the usual Nichols-algebra setting, its Nichols algebra \(\mathcal B(V)\) is the quotient of the braided tensor Hopf algebra \(T(V)\) by its largest homogeneous braided Hopf ideal contained in degrees at least two. For a finitely generated unital \(k\)-algebra \(A\), its Gelfand--Kirillov dimension is
\[
\operatorname{GKdim}_k A
 =\limsup_{m\to\infty}\frac{\log\bigl(\dim_k U^m\bigr)}{\log m},
\]
where \(U\subseteq A\) is a finite-dimensional generating subspace containing \(1\). For an arbitrary unital algebra, one uses the supremum of the Gelfand--Kirillov dimensions of its finitely generated unital subalgebras. These conventions do not supply the missing field or braided data for the source problem.

The source poses the following two requests, originally numbered (3) and (4).
\begin{enumerate}
\item For each \(n\in\{3,4,5\}\), describe more explicitly the geometry underlying the source's reported identification
\[
F(K_n)\simeq T_{!*}(V_n).
\]
The record presents these identifications as facts motivating the question; this draft does not independently verify them. Their precise meaning, including the category in which they hold, requires the missing definitions and surrounding source context.
\item Describe the local systems and sheaves occurring in the intended correspondence for which the corresponding Nichols algebra is known to have finite dimension or finite Gelfand--Kirillov dimension. If the correspondence associates a braided vector space \(V\) over \(k\) to such an object, these conditions mean, respectively,
\[
\dim_k\mathcal B(V)<\infty
\qquad\text{or}\qquad
\operatorname{GKdim}_k\mathcal B(V)<\infty.
\]
The original wording specifically refers to cases where such finiteness is already known. The normalized wording asks more generally about objects whose Nichols algebras satisfy these conditions; the record does not determine whether a complete classification is intended.
\end{enumerate}
These are requests for geometric descriptions, rather than fully specified conjectural propositions. The admissible local systems and sheaves, their base spaces and coefficients, and the required form of a geometric explanation remain unspecified.
\subsection{Short English statement}
Explain geometrically the reported identifications \(F(K_n)\simeq T_{!*}(V_n)\) for \(n=3,4,5\), and describe the local systems and sheaves associated with Nichols algebras of finite dimension or finite Gelfand--Kirillov dimension, retaining the original emphasis on known finite cases. The supplied record omits the definitions and geometric framework needed for a fully precise formulation.
\subsection{Sources}
\begin{itemize}
\item Supplied local UnsolvedMath registry record, permanent ID \(30004836\), source problem code \texttt{catalogue entry}, title \emph{Geometry and Local Systems Behind Finite Nichols Algebras}. The record attributes the questions to the Oberwolfach Reports collection and supplies no individual proposer. Its proposed-year field is \(2021\).
\item Supplied authoritative source citation: \emph{Quantum Groups -- Algebra, Analysis and Category Theory} (2022), Oberwolfach Reports, pp.~2397--2458, DOI \texttt{10.4171/owr/2021/44}. The original excerpt identifies p.~2433 and questions (3) and (4). Source URL and supplied literature reference ``Oberwolfach Report 44'': \url{https://doi.org/10.4171/owr/2021/44}.
\item Supplied literature reference: \emph{Shuffle algebras and perverse sheaves}, \url{https://arxiv.org/abs/1904.09325}. The record supplies no authors or more specific passage for this reference.
\item Literature assessment dated exactly \texttt{2026-08-22}: the supplied curation labels the problem open and states that the sheaf-theoretic correspondence is established in a broad framework, while an explicit geometric explanation of the stated low-rank identifications and a general description of finite-dimensional or finite-Gelfand--Kirillov-dimensional cases remain unknown. This is supplied dated curation derived from the OWR material, not an independently verified literature assessment. The record explicitly calls for independent verification of time-sensitive status and cited claims.
\item The registry labels the statement \texttt{corrected\_verified}, explaining that the questions were isolated from surrounding report prose and their notation normalized. This supplied editorial label does not establish mathematical correctness or independently verify the reported identifications or open status.
\end{itemize}
\end{document}
